\documentclass[11pt]{article}
\usepackage{graphicx}
\usepackage{dcolumn}
\usepackage{amsmath}
\usepackage{multicol}
\usepackage[left=1in,right=1in,top=1in]{geometry}

\newcommand{\nhours}{??}
\newcommand{\nsessions}{??}
\newcommand{\hourspersession}{??}

%\newenvironment{itemize*}%
%  {\vspace{-3.5mm} \begin{itemize}%
%    \setlength{\itemsep}{0pt}%
%    \setlength{\parskip}{0pt}%
%    \setlength{\itemindent}{-1cm}}%
%  {\end{itemize}}

\textwidth 6.5in
\textheight 8.95in
\oddsidemargin -0.1in
%\evensidemargin 0.0in
%\topmargin -0.2truein
\raggedbottom
\pagestyle{empty}

\newcolumntype{d}[1]{D{.}{.}{-#1} }

\newcommand{\Halpha}{H$\alpha$}
\newcommand{\kms}{km s$^{-1}$}
\newcommand{\etal}{et al.}
\newcommand{\msun}{M$_{\odot}$}
\newcommand{\czhelio}{cz$_{\odot}$}
\newcommand{\pmpc}{Mpc$^{-1}$}
\newcommand{\sqd}{deg$^{2}$}
\newcommand{\HI}{H\textsc{i}}

\def\plotfiddle#1#2#3#4#5#6#7{{\centering \leavevmode \vbox to#2{\rule{0pt}{#2}}
     \special{psfile=#1 voffset=#7 hoffset=#6 vscale=#5 hscale=#4 angle=#3}}}

%\pagestyle{empty}
%\setlength{\parindent}{0pt}
%\setlength{\parskip}{6pt}
%\linespread{0.85}

\begin{document}
\parindent 0pt
\baselineskip 12pt
\parskip 8pt
\null
\begin{center}
\vskip -1.2cm
\bf L-Band Wide Observations of Star-forming Galaxies in X-ray Selected Groups: \\
The Quenching of Star Formation in the Group Environment \rm
\end{center}


{\bf Overview:} \hskip 5pt
It is well-established that galaxies in cluster environments have less \HI\ gas (e.g., Solanes {et al} 2011),
%Giovanelli \& Haynes 1985)
lower star formation rates (e.g. Tanaka {et al} 2004), and earlier morphological
types (e.g., Budav\'{a}ri  {et al} 2003) than galaxies in the field.
%Dressler 1980),
These trends appear to hold true also in group environments. Of particular relevance here,
Rasmussen {et al} (2012) have shown that, at fixed
stellar mass, the fraction of group galaxies which are star forming is
lower relative to the comparable fraction in the field.
While physical mechanisms responsible for gas removal
such as ram pressure sweeping and harassment
are understood to be operative in the rich clusters, it is less clear that
these same processes can lead to reduced star formation rates in the lower density environment
of groups.
To constrain these processes requires a robust determination of the threshold in galaxy density,
intracluster medium properties and dark matter halo mass
above which star formation shows evidence of quenching. To this end,
we are undertaking a detailed multiwavelength study of
a set of nearby groups spanning a range of X-ray luminosity, velocity dispersion and richness.
With this proposal, we request time to obtain L-band Wide (LBW) HI observations of star-forming group galaxies
of star forming galaxies which lie below the detection threshhold of the ALFALFA survey. More sensitive
LBW observations will permit us to measure HI masses and star formation efficiences in galaxies
observed to be star-forming but which are
of lower mass and/or gas fraction than those detected by ALFALFA. The study of the full sample
will provide the statistical basis for the determination of the quenching threshold. This proposal requests an
allocation of 122 hours of LBW observations. So that members of the UAT can conduct observations
on-site at Arecibo, we request that some of the time be scheduled to coincide with
the Jan 2014 Undergraduate ALFALFA Team (UAT) workshop and during the ``spring break'' period
in Mar 2014.

{\bf Gas Content and Star Formation:} \hskip 5pt Numerous authors
(e.g. Mahajan {et al} 2010) have presented evidence of the suppression of star
formation in low mass galaxies in both groups and clusters.
Exploring the role of environment on the \HI\ content of galaxies is a
goal of many on-going studies at Arecibo (e.g. Catinella {et al} 2013). For comparative purposes,
the \HI\ masses, gas fractions
(GF; $\equiv$ M$_\text{HI}/$M$_*$), star formation rates
(SFRs), and star formation efficiencies (SFE; $\equiv$ SFR$/$M$_\text{HI}$) have been
quantified for galaxies in the least dense environments and in the
ALFALFA population as a whole (Toribio {et al} 2011; Huang {et al} 2012).
However, investigation of the environmental effects on the links between gas content
and star formation is complicated by known scaling relations: galaxies in denser
environments tend to be higher in stellar mass and
galaxies higher in stellar mass tend to have smaller gas fractions and
less star formation (e.g., Cortese {et al} 2011). These and other
scaling relations between fundamental properties like
stellar mass surface density, star formation rate, and halo spin parameter
(e.g. Huang {et al} 2012), as well as the bias towards higher HI masses
when only detections are included,
must be accounted for in reaching a full
understanding of the mechanisms responsible for altering galaxy evolution
in the group and cluster environment.


Only a few studies
(e.g. Gavazzi {et al} 2013) have attempted to link direct measures of
current star formation to gas content in the same galaxies across a
sufficient range of stellar mass and environment.
To complement the larger statistical
studies being undertaken with the full ALFALFA dataset, we are
investigating the detailed distribution and properties of the
population of galaxies {\it known to be star-forming}
in a sample of groups (see Table 1) in the sky area covered by ALFALFA
but of varying
richness, morphological type mix, X-ray luminosity, and velocity
dispersion. We use not only the \HI\ measures from ALFALFA but also
complementary data available in the literature and in digital
databases such as the NASA-Sloan Atlas (NSA; Blanton {et al}
2011). Additionally, in order to explore specifically the quenching of
star formation on the scale of these groups, we have obtained
wide-field ($\ge$ 1 \sqd) \Halpha\ and broadband R images using the
MOSAIC camera on the
KPNO 0.9m telescope; multiple offset fields were observed to give
radial extent.  The wide-field
\Halpha\ imaging allows us to establish the rate of
on-going massive star formation as well as its distribution in group
member galaxies. Still, because the \Halpha\ does not generally cover
the whole group, we supplement the sample with additional star forming
galaxies chosen on their basis of line ratios extracted from the
NSA following the diagnostic scheme of Kewley {et al} (2006).
With this wealth of optical data on our groups we can use the
\Halpha\ and UV
fluxes in combination with \HI\ measures to determine the SFEs. However,
at the distances of many of these groups,
the lower mass galaxies as well as gas-depleted ones,
are expected to have \HI\ fluxes below ALFALFA's
well-behaved sensitivity limits (Haynes {et al} 2011) so that ALFALFA
measures alone are insufficient to establish the GFs and SFEs for many
galaxies which have known \Halpha\ emission, but are not yet detected
in \HI\. To rectify this
limitation, we propose L-band wide (LBW) observations to obtain
measures of \HI\ gas content of a complete set of galaxies
known to be star forming from the optical data but which were {\it not}
detected by ALFALFA, thereby allowing us to measure the GF and SFE to
lower \HI\ masses, into the regime of gas deficiency. The increased
sensitivity is critical to establish the threshold of gas depletion, star
formation quenching, and morphological transformation within these
intermediate density environments and their surrounding extended
supercluster substructures. The results of the statistical study of the 12 groups
will inform us on the choice of selected fields which we intend
to propose to map in \HI\ in greater detail with the VLA to see how the
distribution of the \HI\ and \Halpha\ correlates in individual galaxies
close to the quenching threshold.

\begin{center}
\begin{figure}[t]
\centering
\includegraphics[scale=0.25]{n5178R.eps}
\hspace{0.1in}
\includegraphics[scale=0.25]{n5178Ha.eps}
\caption{NGC 5178 imaged at R (left) and \Halpha\ (right) wavelengths
at the WIYN 0.9m telescope at KPNO.
The galaxy is classified as S0 and is located in the MKW11 group. It
has clear star formation, but it is not detected by ALFALFA. The star
formation is truncated, reminiscent of environmentally altered
Virgo Cluster galaxies in Koopmann \& Kenney (2004).
\label{haimage}}
\end{figure}
\end{center}
\vskip -1.5cm

{\bf Sample Selection:} \hskip 5pt The target group sample (Table 1) has been selected to meet the
following criteria: (i) in the ALFALFA-SDSS overlap region available to us in January 1, 2013;
(ii) cz $<$ 8500 km/s (Dist$<$ 120Mpc); (iii) surveyed at
X-ray wavelengths; (iv) excluding rich clusters such as Coma; and (v) observed by us in
\Halpha\ in April 2013. The target list of star-forming galaxies has been drawn up by searching
the NSA, the ALFALFA catalog and the list of \Halpha\ detections in our
images and which lie within 2 Mpc of the group center and 5$\sigma$
of its systemic velocity. The LBW target list consists of those galaxies which are
{\it not} detected by ALFALFA but which {\it are classified} as star-forming either by
the detection of H$\alpha$ emission in our KPNO images or
based on the NSA emission line flux ratio "star forming" classification according
to Kewley {et al} (2006).
The broadband R and narrowband H$\alpha$ image of a typical H$\alpha$ {\it but
not} ALFALFA detection is shown in Figure 1. In Table 1,
we include the principal properties of each group as well as the number of galaxies
including in the various subsets: detected by ALFALFA, detected in the H$\alpha$
images but not by ALFALFA, exhibiting central star-forming characteristics but
not detected by ALFALFA and the total number of LBW targets.

{\bf Observing Strategy:} \hskip 5pt The proposed LBW pointed observations consist of 5 minute ON-OFF total
power pairs with the interim correlator covering a 25 MHz bandwidth
with 9-level sampling. We use doppler tracking and center the spectra
at the frequency appropriate to each target's redshift. In order to optimize the observing
schedule, we will adopt a flexible multiple-pass strategy designed to minimize slew time
and to integrate deeper only when needed.
Each galaxy will be observed once during a first night, with
subsequent nights used to perform additional ON-OFF observations of galaxies which
remain undetected in \HI. We expect many galaxies to be detectable with a single
ON-OFF pair, which, assuming a system temperature of 27K and a velocity resolution of
10 km/s, should attain an rms value of $\sim0.65$ mJy. A maximum of 3 ON-OFF pairs
gives an rms of $\sim0.37$ mJy which, assuming a detection limit S/N of $\sim5$
translates to an estimated \HI\ mass limit of
$10^{7.4}$ M$_\odot$ in the
closest group ($\sim25$ Mpc) and $10^{8.7}$ M$_\odot$ in the furthest ($\sim100$ Mpc).
These observations would represent up to a factor of 6 improvement in rms flux over
ALFALFA observations. Physically this means that while ALFALFA could detect a galaxy of
normal HI content (GF $\sim 2$ as given by Eqn. 1 of Huang {et al} 2012) of stellar mass
$10^9$ M$_\odot$ at 100 Mpc, these observations will detect such galaxies
with gas fractions of as low as 1/3.
On average, we expect to require two ON-OFF pairs
per galaxy, which, with slew and calibration time, require slightly less than 30 minutes
per galaxy.


\bf Observing Request: \rm Our target list includes 235 sources within 12 groups. Given the
practical rate of a bit more than 2 sources/hour plus startup/end overhead, we request an allocation
of 122 hours with LBW and the interim correlator. The UAT will be heavily involved in all aspects
of this program; we would like to conduct a significant fraction of the
observations at Arecibo. Considering the sky distribution of the
target fields and in order to allow UAT faculty and students to participate in on-site observing,
our request of 122 hours is distributed as follows:
\begin{center}
\vskip -8pt
\begin{tabular}{rcl}
3 sessions  & 08$^h <$ LST $< 14^h$             & coincident with the UAT workshop (Jan 12-15, 2014)\\
4 sessions  & 08$^h30^m <$ LST $< 16^h30^m$      & in late March 2014\\
12 sessions & 10$^h30^m <$ LST $< 16^h30^m$      & some of which would be in late March 2014
\end{tabular}
\end{center}

{\bf ALFALFA and LBW Status/Productivity Report}:\hskip 5pt
The legacy ALFALFA survey drift scan observations were completed in October 2012; all
data have been Level I processed; gridding, source extraction, catalog and data product
production continues
with a next significant data release expected in the 2nd quarter of 2014. The paper associated with the
next data release will include the statistical analysis of the low signal-to-noise ALFALFA
sources which have been the targets of the LBW programs conducted in 2012 and spring 2013
(A2669, A2707 and A2752).

{\bf Broader Impact:}\hskip 5pt The UAT is a consortium of
19 undergraduate-focused institutions that provides undergraduate research
experiences centered on the ALFALFA survey to their students. 183 undergraduates have
participated during the 6 years of the program, including 106 students who
attended workshops at Arecibo and 47 who have observed onsite for ALFALFA
and LBW runs. Multiple UAT institutions
collaborate on a common project (the UAT Groups Project) to quantify the impact of environment
on galaxies in groups; this proposal is directly related to that project. Team members
identify analysis methods and needed software and commit
to producing and sharing software tools and analysis results.
Students have written and documented many of the software tools, and a number
have been incorporated in the general ALFALFA-IDL package distributed to the larger
ALFALFA collaboration. In addition to the undergraduate involvement,
three Cornell University Ph.D. theses were completed in summer 2013: E. Adams (now at ASTRON),
S. Huang (now at ASIAA), E. Papastergis (now at U. Groningen). Three continuing Cornell
graduate students are co-Is on the current proposal because of their involvement with
UAT activities.


{\bf ALFALFA-Based Refereed Publications (Since Mar 1, 2013 Only):}\\
\parskip 0pt
\parindent 0pt
$\bullet$ \hskip 3pt {\it A Catalog of Ultra-compact High Velocity Clouds from the ALFALFA Survey: Local Group Galaxy Candidates?},
Adams, E.A.K.$^\dagger$, Giovanelli, R. \& Haynes, M.P. 2013, ApJ 768, 77.

$\bullet$ \hskip 3pt {\it ALFALFA Discovery of the Nearby Gas-Rich Dwarf Galaxy Leo P: \HI\ Properties},
Giovanelli, R., Haynes, M.P., Adams, E.A.K.$^\dagger$, Cannon, J.M., Rhode, K.L.,
Salzer, J.J., Skillman, E.D., Berstein-Cooper, E.Z.$^*$, \& McQuinn, K.B. 2013, AJ 146, 15

$\bullet$ \hskip 3pt {\it ALFALFA Discovery of the Nearby Gas-Rich Dwarf Galaxy Leo P:
Distance Estimate and Results from Optical Imaging of Leo P},
Rhode, K.L., Salzer, J.J., Haurberg, N.C.$^\dagger$,
Van Sistine, A.$^\dagger$, Young, M.D.$^\dagger$, Haynes, M.P., Giovanelli, R.,
Cannon, J.M., Skillman, E.D., McQuinn, K.B.W., \& Adams, E.A.K.$^\dagger$ 2013, AJ 145, 149

$\bullet$ \hskip 3pt {\it ALFALFA Discovery of the Nearby Gas-Rich Dwarf Galaxy Leo P: An Extremely Metal Deficient Galaxy},
Skillman, E.D., Salzer, J.J., Berg, D.A.$^\dagger$, Pogge, R.W., Haurberg, N.C.$^\dagger$, Cannon, J.M.,
Aver, E., Olive, K.A. Giovanelli, R., Haynes, M.P., Adams, E.A.K.$^\dagger$, McQuinn, K.B.W. \& Rhode, K.L.
2013, AJ 146, 3

$\bullet$ \hskip 3pt {\it The Clustering of ALFALFA Galaxies: Dependence on \HI\ Mass, Relationship to Optical Samples
and Clues on Host Halo Properties},
Papastergis, E.$^\dagger$, Giovanelli, R., Haynes,
   M.P., Rodriguez-Puebla, A. \& Jones, M.G.$^\dagger$ 2013 Ap.J (in press)

$\bullet$ \hskip 3pt {\it ALFALFA Discovery of the Nearby Gas-Rich Dwarf Galaxy Leo P: IV. Distance Measurement from LBT Optical Imaging},
McQuinn, K.B.W., Skillman, E.D., Berg, D.$^\dagger$ , Cannon, J.M., Adams,
E.A.K.$^\dagger$ , Dolphin, A., Giovanelli, R., Haynes, M.P., Rhode, K.L. \&
Salzer, J.J. 2013 (submitted)

\vskip 5pt {\bf $^\dagger$ Graduate student; $^*$ Undergraduate student}

%\newpage
\vskip 12pt
\begin{table}[h!]
\begin{center}
\small
\caption{Properties of the UAT Group Sample}
\begin{tabular}{lccccccccd{2}}
\hline
 \small{Group} & \small{RA} & \small{Dec} & \small{$v_{\odot}$} & \small{$\sigma_v$} & \small{N$_\text{HI}$} & \small{N$_{\text{H}\alpha}$} & \small{N$_\text{SF}$} & \small{N$_\text{LBW}$} & \multicolumn{1}{c}{{\small log L$_\text{x}$}} \\
  &  & & \small{(\kms)} & \small{(\kms)} & &  & & & \multicolumn{1}{c}{\small (ergs s$^{-1}$)}\\
	& \small (1) & \small (2) & \small (3) & \small (4) & \small (5) & \small (6) & \small (7) & \small (8) &  \multicolumn{1}{c}{\small (9)} \\
	\hline
NRGb041 &  09 25 12.2 & +11 16 49 &   3616 & 208 & 10 & 2 & 3 &	5   &<41.50  \\
NRGb054 &  09 41 16.7 & +11 33 16 &   6460 & 331 & 11 & 10 & 11 & 21  & <41.80 \\
NRGb151 &  11 42 09.4 & +10 16 30 &   6265 & 257 & 13 & 1 & 4  & 5   &  42.04 \\
NRGb157 &  11 47 51.8 & +13 16 33 &   3745 & 457 & 31 & 1 & 6  & 7   &<41.50 \\
NRGb168 &  11 58 30.5 & +25 08 23 &   4789 & 162 &  8 & 7 & 5  & 12  &<41.60  \\
NRGb206 &  12 24 02.6 & +09 21 25 &   7256 & 457 & 21 & 8 & 19 & 27  &<41.90 \\
NRGb247 &  13 29 31.2 & +11 47 19 &   7056 & 269 & 22 & 15& 14 & 29  &  42.67 \\
NRGb282 &  14 05 57.0 & +09 07 58 &   6872 & 660 & 22 & 33& 8  & 41  &<42.00 \\
NRGb301 &  14 28 01.9 & +25 34 10 &   4819 & 302 & 24 & 9 & 12 & 22  & <41.70 \\
MKW 8   &  14 40 49.1 & +03 24 11 &   8100 & 325 & 10 & 21& 18 & 39  & 41.98 \\
NGC 5846&  15 05 32.0 & +01 41 48 &   1749 & 322 &  7 & 9 & 13 & 22  &  42.00  \\
NRGb331 & 15 06 41.4 & +12 44 13 &   6904 & 204 &  9 & 2 & 3  & 5  & <42.20 \\
\hline
\end{tabular}
\end{center}
\small
{\bf Notes:} \hskip 5pt
Columns 1 and 2: Positions of cluster center;
Columns 3 and 4: Heliocentric cluster central velocity and velocity dispersions;
Column 5: Number of group galaxies detected in \HI\ by the ALFALFA survey;
Column 6: Number detected in \Halpha\ imaging but not ALFALFA (numbers in parentheses are preliminary); Column 7: Number classified by emission line ratios in the NASA-Sloan Atlas as star-forming,
but not detected in ALFALFA or \Halpha\ images;
Column 8: Number of group galaxies we propose to observe, the sum of columns 6 and 7;
Column 9: X-ray luminosities. Properties in Columns 1-4 and 9 are from
Mahdavi {et al} (2000), except for NGC 5846 which has X-ray
luminosity from Boehringer {et al} (2000) and velocity and dispersion
from Mahdavi, Trentham, \& Tully (2005).
%N$_{H\alpha}$: Number of group galaxies detected in \Halpha\ images but not in ALFALFA\\
%N$_{H\alpha}$: Number of group galaxies in the NASA-Sloan Atlas and classified by
%their line ratios as star-forming but not detected in ALFALFA or in our \Halpha\ images
%(Note that our \Halpha\ images do not generally span the full $4\times4$ Mpc region sampled by
%the NASA-Sloan Atlas.) \\
%Positions (X-ray centroids where detected), central velocities
%($v_{he}$) , velocity dispersions ($\sigma_v$), and X-ray luminosities
%are from Mahdavi {et al} (2000) except for NGC 5846, which has X-ray
%luminosity from Boehringer {et al} (2000) and velocity and dispersion
%from Mahdavi, Trentham, \& Tully (2005)}
\end{table}



\begin{center}
\bf References:\\
\end{center}
\begin{multicols}{2}
%\parskip 0pt
%\parindnt 0pt
\small

%\noindent Balogh, M. L.  et al. 1997, ApJ, 488, 75

\noindent Blanton, M. et al. 2011, AJ, 142, 31

\noindent Boehringer, H. et al. 2000, ApJS, 129, 435

\noindent Budav\'{a}ri, T. et al. 2003, ApJ 595, 59

\noindent Catinella, B. et al. 2013, MNRAS (in press); astro-ph 13.08,1676

%\noindent Chung, A. et al. 2007, ApJ, 659, L1

\noindent Cortese, L. et al. 2011, MNRAS, 415, 1797

%\noindent Cortese, L. \& Hughes, T.M. 2009, MNRAS 400, 1225

%\noindent Dale, D. A. et al. 2001, AJ, 121, 1886

%\noindent Fabello, S. et al. 2012, MNRAS, 427, 2841

%\noindent Fosatti, M. et al. 2013, A\&A, 553, 91

\noindent Gavazzi. G. et al. 2013, A\&A, 553, 89

%\noindent Giovanelli, R. \& Haynes M. P. 1983, AJ, 88, 881

%\noindent Giovanelli, R. et al. 2005, AJ, 130, 2598

%\noindent Gomez, P. L., et al. 2003, ApJ, 584, 210

%\noindent Gunn, J. E. \& Gott, J. R, ApJ, 176, 1

\noindent Haynes, M.P., et al. 2011, AJ 142, 170

\noindent Huang, S., et al. 2012, ApJ. 756, 113

%\noindent Johnson, K. E. et al. 2007, AJ, 134, 1522

%\noindent Kilborn, V.A. et al. 2009, MNRAS, 400, 1962

\noindent Kewley, L.J. et al. 2006, MNRAS, 372, 961
%Groves, B., Kauffmann, G., \& Heckman, T.

\noindent Koopmann, R.A. \& Kenney, J.D.P. 2004, ApJ, 613, 851

%\noindent Koopmann, R. A. \& Kenney, J. D. P. 2006, ApJs, 162, 97

%\noindent Koopmann, R. A., Kenney, J. D. P., \& Young 2001, ApJs, 135, 125

%\noindent Larson, R. B., Tinsley, B. M., Caldwell, C. N. 1980, ApJ, 237, 692

%\noindent Lewis, I., et al. 2002, MNRAS, 334, 673

%\noindent Lu, T. et al.  2012, MNRAS, 24, 126

\noindent Mahajan, S. et al.  2010, MNRAS 404, 1745

\noindent Mahdavi, A. et al. 2000, ApJ, 534, 114

\noindent Mahdavi, A., Trentham, N., Tully, R.B. 2005, AJ, 130, 1502

%\noindent Moore, B. et al. 1996, Nature, 379, 613

%\noindent Poggianti, B. et al. 1999, ApJ, 518, 576

%\noindent Postman, M. et al. 2005, ApJ, 623, 721

\noindent Rasmussen, J. et al. 2012, ApJ 757, 122

%\noindent Rose, J.A., Robertson, P., Miner, J., \& Levy, L. 2010, AJ, 139, 765

%\noindent Sakai, S., Kennicutt, R. C., van der Hulst, J. M., \& Moss, Chris 2002, ApJ, 578, 842

\noindent Solanes, J.M. et al. 2001, ApJ, 548, 97

\noindent Tanaka, M. et al. 2004, AJ 128, 2677

\noindent Toribio, M.C., et al.  2011, ApJ 732, 92\\

%\noindent Tzanavaris, P. et al 2010, ApJ, 716, 556

%\noindent Verdes-Montenegro, L., et al.  2001, A\&A, 377, 812

%\noindent York, D. G. 2000, AJ, 120, 1579

\end{multicols}

\end{document}
